\section{Projection and tangent positivity}
\label{sec:tangent}

The sequence theorem will be proved by two inductions: first for
all-high data, and then for data with a low factor. In a counterexample
of least ambient dimension, quotienting by a product of abelian
subvarieties must fail to preserve the strict codimension inequality.
We express this obstruction precisely and use it to prove two
positivity statements for every subproduct through the marked points.
The first statement will feed the height inequality; the second will
distinguish measures in the bounded-height argument.

\subsection{Quotients of minimal counterexamples}

\begin{lemma}[Projection exclusion]
\label{lemma:projection}
Suppose that sequence data as in Definition~\ref{def:sequence-data},
together with $V$ and $\Phi$ as in Theorem~\ref{thm:sequence}, satisfy
\eqref{eq:forbidden-sequence} and have least $\dim\cA$ among
counterexamples with the same presence or absence of a low factor.
If a low factor occurs, assume also that Theorem~\ref{thm:sequence}
has already been proved for all-high data in every dimension.
There do not exist fixed abelian subvarieties
$C_\ell\subseteq A_\ell$, with $C=\prod_\ell C_\ell\ne0$,
and nonnegative integers $a_\ell$ such that, on an infinite
subsequence,
\begin{equation}
 \label{eq:projection-obstruction}
 \dim_{y_{i\ell}}\bigl(Y_\ell\cap(y_{i\ell}+C_\ell)\bigr)
 \ge a_\ell\quad(\ell\in I),
 \qquad
 \sum_\ell a_\ell\ge\rank(\Phi|_{\Lie(C)}).
\end{equation}
\end{lemma}

\begin{proof}
Suppose such subvarieties and integers exist and restrict to the given
subsequence. Let $\pi_\ell:A_\ell\to A'_\ell=A_\ell/C_\ell$
be the quotient, let $W'_\ell=\pi_\ell(W_\ell)$, and let
$e_\ell$ be the dimension of the image of $C_\ell$ in
$D_\ell=A_\ell/W_\ell$. We first construct fixed irreducible
parents $Y'_\ell\subseteq A'_\ell$ for the projected points with
local $W'_\ell$-fiber bounds
\begin{equation}
 \label{eq:projected-fiber-bound}
 m'_\ell=m_\ell+e_\ell-a_\ell.
\end{equation}

Fix one factor and temporarily omit its index. In $Y$ let $S$ be the
intersection of the open locus where local fiber dimension over $A/W$
is at most $m$ and the closed locus where local fiber dimension over
$A/C$ is at least $a$. Local fiber dimension is upper semicontinuous,
so $S$ is locally closed and contains every marked point under
consideration. Every nonempty fiber of $S\to A/C$ has dimension at
least $a$. Indeed, through each of its points there is an irreducible
component of the original $C$-fiber of dimension at least $a$. That
whole component belongs to the closed high-fiber locus, and its
intersection with the open low-fiber locus is a nonempty open subset
of the same dimension.

Put $E=A/(W+C)=A'/W'$. The map from a fiber of $S\to E$ to
$A/W$ has image in an $e$-dimensional coset, and all of its nonempty
fibers have dimension at most $m$. Consequently every fiber of
$S\to E$ has dimension at most $m+e$. Let $Q$ be the constructible
image of $S$ in $A'$. The fibers of $S\to Q$ have dimension at least
$a$, so the dimension inequality for fibers, applied over every
point $t\in E$, gives
\[
 \dim Q_t\le m+e-a.
\]
This can also be seen by applying the fiber dimension theorem to
each irreducible locally closed stratum of $Q_t$: one of the finitely
many strata in its inverse image dominates it with generic fiber
dimension at least $a$. In particular $m+e-a\ge0$ when a marked
point exists.

Partition $Q$ into finitely many irreducible locally closed subsets.
After taking a subsequence, the projected points lie in one such
subset $Q^\circ$. Set $Y'=\overline{Q^\circ}$. Because
$Q^\circ$ is open in its closure, its local germ at any of the
selected points equals the germ of $Y'$. The local fiber dimension
of $Y'\to E$ at that point is therefore at most $m+e-a$. Performing
this construction in each of the finitely many factors, with
successive subsequences, proves \eqref{eq:projected-fiber-bound}.
All these parents are defined over $\Qbar$.

The new parents provide the local fiber bounds. We next check the
small-projection condition and the strict codimension inequality.
Let $\pi=\prod_\ell\pi_\ell$ and
$V'=\pi(V)\subseteq\Lie(W')$. This is the image of a real
homomorphism, so Lemma~\ref{lemma:real-calculus} supplies its
orthogonal projector. If $\Phi'$ is the complementary projector,
then
\[
 \Phi'\pi=\Phi'\pi\Phi,
\]
since $\pi(V)=V'$. Thus the projected normalized points satisfy
the required small-projection condition. The height upper bounds
are preserved by the fixed quotient homomorphisms.

The gain in fiber bounds compensates for the rank lost under projection:
\begin{align}
 \codim_{\Lie(W')}V'
 &=\dim W-\dim C+\sum_\ell e_\ell
      -\dim V+\dim(V\cap\Lie(C))\notag\\
 &=\codim_{\Lie(W)}V+\sum_\ell e_\ell
      -\rank(\Phi|_{\Lie(C)})\notag\\
 &>\sum_\ell(m_\ell+e_\ell-a_\ell)
   =\sum_\ell m'_\ell.
 \label{eq:projection-codimension}
\end{align}
Here $\dim(W\cap C)=\dim C-\sum e_\ell$ and
$\ker(\Phi|_{\Lie(C)})=V\cap\Lie(C)$.

It remains to obtain nonzero distinguished factors and check the
scale conditions. Discard each factor with $W'_\ell=0$. In such a
factor $V'$ has zero coordinate, so deletion does not change the codimension inside
$\Lie(W')$ and only decreases the sum of the nonnegative bounds
$m'_\ell$. The small-projection condition on the remaining product
follows by composition with its coordinate projection. At least one
factor survives, since otherwise
\eqref{eq:projection-codimension} would say that zero is strictly
greater than a sum of nonnegative integers. The scales on the
remaining ordered list still have successive ratios tending to zero.
For a surviving high factor, a homomorphism nonzero on $W'_\ell$
composes with $\pi_\ell$ to one nonzero on $W_\ell$, proving
\eqref{eq:sequence-nondegenerate}. If the low factor survives, then
$W'_0=A'_0$ and its torsion-coset avoidance follows by pulling a
proper torsion coset back under $\pi_0$.

We have obtained counterexample data of smaller total ambient
dimension because $C\ne0$. If the low factor is present both before
and after projection, or absent both before and after, this
contradicts minimality. If it has disappeared, it contradicts the
already-established all-high case. This proves the exclusion.
\end{proof}

\subsection{A product tangent obstruction}

Let $Z=\prod_{\ell\in I}Z_\ell$ be a product of subvarieties of
$\prod_\ell A_\ell$, and let $g$ be a complex linear map from its
ambient Lie algebra to a Hermitian vector space with positive
$(1,1)$-form $\omega$. Define
\begin{equation}
 \label{eq:tangent-integral}
 P_Z(g)=\int_Z(g^*\omega)^{\dim Z},
\end{equation}
where the pullback is regarded as an invariant form on the ambient
abelian variety. We use the usual integral over the smooth locus,
or equivalently the integral on a resolution. When polarizations
are involved, the invariant forms represent their first Chern
classes. If $\dim Z=0$, the integral is the degree of the reduced
point, hence is one. In every dimension semipositivity gives
\[
 P_Z(g)>0\quad\Longleftrightarrow\quad
 g\text{ is injective on }T_zZ
 \text{ for some smooth }z\in Z.
\]
For positive scalars $b_\ell$ the direct-sum decomposition of a
product tangent space gives
\begin{equation}
 \label{eq:product-scaling}
 P_Z\bigl(g\operatorname{diag}(b_\ell)\bigr)
 =\left(\prod_\ell b_\ell^{2\dim Z_\ell}\right)P_Z(g).
\end{equation}
Indeed each factor scaling acts on the complex determinant of its
tangent space by $b_\ell^{\dim Z_\ell}$, and the associated real
top form acquires the square of its absolute value. The same formula
holds for nonzero complex scalars with $|b_\ell|$ in place of
$b_\ell$.

For a subvariety $D$ of a complex abelian variety, its
\emph{geodesic closure} $\langle D\rangle_{\mathrm{geo}}$ is the
smallest coset of an abelian subvariety containing $D$. The translating
point need not be torsion. We will use Ax's theorem to turn failure of
tangent injectivity into a product direction having large fibers.

\begin{lemma}[Product tangent obstruction]\label{lem:ax-product}
Let $A_\ell$, $\ell\in I$, be complex abelian varieties, let
$Z_\ell\subseteq A_\ell$ be subvarieties, and put
$Z=\prod_{\ell\in I}Z_\ell$ and $A=\prod_{\ell\in I}A_\ell$.
Let $g\colon\Lie(A)\longrightarrow E$ be a complex linear map to a
Hermitian vector space. If $P_Z(g)=0$, there are abelian subvarieties
$C_\ell\subseteq A_\ell$ such that, for every
$z=(z_\ell)_\ell\in Z$,
\begin{equation}\label{eq:ax-product-obstruction}
 \sum_{\ell\in I}
 \dim_{z_\ell}\bigl(Z_\ell\cap(z_\ell+C_\ell)\bigr)
 >\rank\left(g\bigm|_{\bigoplus_{\ell\in I}\Lie(C_\ell)}\right).
\end{equation}
In particular, $\prod_\ell C_\ell$ is nonzero.
\end{lemma}

\begin{proof}
We will replace the subgroup supplied by Ax's theorem by a product
subgroup. The degeneration used for this replacement cannot decrease
the tangent-intersection dimension. Its $g$-rank must also be bounded
by the rank appearing in Ax's inequality. We arrange this comparison
first by choosing factor scalings that maximize the rank on every
abelian subgroup; only then do we apply Ax's theorem.

Write $D(\lambda)=\operatorname{diag}(\lambda_\ell)$ for
$\lambda\in(\mathbb C^\times)^I$. For every abelian subvariety
$B\subseteq A$, the rank of $gD(\lambda)|_{\Lie(B)}$ attains its maximum
on a nonempty Zariski open subset of $(\mathbb C^\times)^I$.
There are only countably many abelian subvarieties of $A$: their rational
homology groups determine distinct rational subspaces of $H_1(A,\mathbb Q)$.
The Baire category theorem therefore gives one $\lambda$ at which all
these ranks are maximal. Fix such a $\lambda$ and set $h=gD(\lambda)$.

At a smooth point of $Z$, its tangent space is the direct sum of the
tangent spaces of its factors. It is consequently preserved by
$D(\lambda)$. The failure of tangent injectivity for $g$, which is
equivalent to $P_Z(g)=0$, thus also holds for $h$ at every smooth point.
This is the same product property underlying
\eqref{eq:product-scaling}.

Choose a very general smooth point $z\in Z$. We require that the rank of
$h$ on $T_zZ$ be maximal, and that, for every abelian product subgroup
$C\subseteq A$, the differential of $Z\longrightarrow A/C$ have its
generic rank at $z$. These requirements exclude only countably many
proper closed subsets of $Z$, so they can be imposed simultaneously.
On a sufficiently small logarithmic chart around $z$, the map obtained
by applying $h$ to the logarithm has constant rank less than $\dim Z$.
Its fiber through $z$ therefore has a positive-dimensional irreducible
analytic germ. Let $J$ be a sufficiently small irreducible representative
in the logarithmic chart, let $U_1$ be the Zariski closure of $\exp(J)$,
and let $z+C^*$ be the smallest coset containing $U_1$.
Writing $\widetilde z$ for the logarithm of $z$, we have
\[
 J\subseteq\widetilde z+
 \bigl(\Lie(C^*)\cap\ker h\bigr).
\]

We use Ax's theorem \cite{Ax1972}, in the precise analytic-subset formulation
\cite[Theorem~5.3]{HabeggerPila2016}. If an irreducible analytic subset
$K$ of an open subset of a translate of a complex vector space $U$
has exponential image with Zariski closure $B$, that theorem gives
\[
 \dim\langle B\rangle_{\mathrm{geo}}-\dim B
 \leq\dim U-\dim K,
\]
where $\langle B\rangle_{\mathrm{geo}}$ denotes the smallest abelian
coset containing $B$. Applied to $J$ and
$U=\Lie(C^*)\cap\ker h$, it gives
\[
 \dim U_1\geq
 \rank\bigl(h|_{\Lie(C^*)}\bigr)+\dim J
 >\rank\bigl(h|_{\Lie(C^*)}\bigr).
\]
Since $U_1\subseteq Z\cap(z+C^*)$ contains $z$, the Zariski tangent
space at $z$ gives
\begin{equation}\label{eq:ax-tangent-before-degeneration}
 \dim\bigl(T_zZ\cap\Lie(C^*)\bigr)
 \geq\dim U_1
 >\rank\bigl(h|_{\Lie(C^*)}\bigr).
\end{equation}
No smoothness assumption on $U_1$ at $z$ is needed for this inequality.

We now replace $C^*$ by a product direction. Order the index set $I$
and choose strictly increasing integers $w_\ell$. Put
$S=\Lie(C^*)$ and $F_\ell=\bigoplus_{j\geq\ell}\Lie(A_j)$.
A basis of $S$ adapted to the filtration $S\cap F_\ell$ shows that
the Grassmannian limit
\[
 S_0=\lim_{t\longrightarrow0}
       \operatorname{diag}(t^{w_\ell})S
     =\bigoplus_{\ell\in I}\operatorname{pr}_\ell(S\cap F_\ell)
\]
exists. Indeed, after dividing an adapted basis vector by the power of
$t$ corresponding to its first nonzero coordinate, its limit is the
projection to that coordinate. Each summand is the Lie algebra of an
abelian subvariety $C_\ell\subseteq A_\ell$: take the identity component
of the intersection of $C^*$ with the product subgroup corresponding
to $F_\ell$, and project to $A_\ell$.

The product tangent space $T_zZ$ is preserved by all the displayed
diagonal scalings. Upper semicontinuity of intersection dimension on
the Grassmannian therefore gives
\[
 \dim(T_zZ\cap S_0)\geq\dim(T_zZ\cap S).
\]
Set $R=\rank(h|_S)$. By the choice of $\lambda$, for every $t\neq0$,
\[
 \rank\left(g\bigm|_{\operatorname{diag}(t^{w_\ell})S}\right)
 =\rank\bigl(g\operatorname{diag}(t^{w_\ell})|_S\bigr)\leq R.
\]
The condition that the restriction of a fixed linear map have rank
at most $R$ is closed on the Grassmannian. Hence
$\rank(g|_{S_0})\leq R$. Together with
\eqref{eq:ax-tangent-before-degeneration}, this proves
\[
 \dim\left(T_zZ\cap\bigoplus_\ell\Lie(C_\ell)\right)
 >\rank\left(g\bigm|_{\bigoplus_\ell\Lie(C_\ell)}\right).
\]

By the choice of $z$, the dimension on the left is the generic fiber
dimension of $Z\longrightarrow A/\prod_\ell C_\ell$.
Local fiber dimension at every point is at least this generic value.
Since both $Z$ and the subgroup are products, local fiber dimensions
add over the factors. This proves
\eqref{eq:ax-product-obstruction} for every point of $Z$.
\end{proof}

\subsection{Finite directions in fixed parents}

We next explain why the subgroup directions obtained from this lemma
can be chosen from finite sets when the varieties $Z_\ell$ vary inside
fixed parents. For a subvariety $D$ of an abelian variety, put
\[
 \delta_{\mathrm{geo}}(D)
 =\dim\langle D\rangle_{\mathrm{geo}}-\dim D.
\]
A subvariety $D\subseteq Y$ is \emph{geodesic-optimal in $Y$} if every
proper enlargement $D\subsetneq E\subseteq Y$ has
$\delta_{\mathrm{geo}}(E)>\delta_{\mathrm{geo}}(D)$.

\begin{lemma}[Finite directions for local rank inequalities]
\label{lem:finite-directions}
Fix subvarieties $Y_\ell\subseteq A_\ell$ of complex abelian varieties.
For each $\ell$, there is a finite set $\mathcal H_\ell$ of abelian
subvarieties of $A_\ell$ with the following property.
Suppose that $y_\ell\in Z_\ell\subseteq Y_\ell$, that
$C_\ell\subseteq A_\ell$ are abelian subvarieties, and that
$g\colon\bigoplus_\ell\Lie(A_\ell)\longrightarrow E$ is complex linear.
Put
\[
 d_\ell=\dim_{y_\ell}\bigl(Z_\ell\cap(y_\ell+C_\ell)\bigr).
\]
If $\varepsilon\in\{0,1\}$ and
\[
 \sum_\ell d_\ell\geq
 \rank\left(g\bigm|_{\bigoplus_\ell\Lie(C_\ell)}\right)
 +\varepsilon,
\]
there are $H_\ell\in\mathcal H_\ell$ such that
\begin{equation}\label{eq:fixed-parent-rank}
 \sum_\ell\dim_{y_\ell}
 \bigl(Y_\ell\cap(y_\ell+H_\ell)\bigr)
 \geq\rank\left(g\bigm|_{\bigoplus_\ell\Lie(H_\ell)}\right)
 +\varepsilon.
\end{equation}
If $\sum_\ell d_\ell>0$, the product $\prod_\ell H_\ell$ is nonzero.
\end{lemma}

\begin{proof}
Choose an irreducible component $D_\ell$ of
$Z_\ell\cap(y_\ell+C_\ell)$ through $y_\ell$ of dimension $d_\ell$.
Write its smallest containing coset as $y_\ell+B_\ell$; then
$B_\ell\subseteq C_\ell$. Among the subvarieties of $Y_\ell$
containing $D_\ell$ and having geodesic defect at most
$\delta_{\mathrm{geo}}(D_\ell)$, choose one, denoted $E_\ell$, of
maximum dimension. Such a choice is possible because $D_\ell$ is an
admissible choice and dimensions are bounded integers. It is
geodesic-optimal: a proper enlargement of no greater defect would
contradict its choice.

Write $\langle E_\ell\rangle_{\mathrm{geo}}=y_\ell+H_\ell$.
The inclusion $D_\ell\subseteq E_\ell$ implies
$B_\ell\subseteq H_\ell$. The defect comparison gives
\[
 \dim H_\ell-\dim B_\ell
 \leq\dim E_\ell-d_\ell.
\]
Rank increases by at most the dimension added to the domain. Therefore
\begin{align*}
 \rank\left(g\bigm|_{\bigoplus_\ell\Lie(H_\ell)}\right)
 &\leq
 \rank\left(g\bigm|_{\bigoplus_\ell\Lie(B_\ell)}\right)
 +\sum_\ell(\dim H_\ell-\dim B_\ell)\\
 &\leq
 \rank\left(g\bigm|_{\bigoplus_\ell\Lie(C_\ell)}\right)
 +\sum_\ell(\dim E_\ell-d_\ell)\\
 &\leq\sum_\ell\dim E_\ell-\varepsilon.
\end{align*}
Since $E_\ell\subseteq Y_\ell\cap(y_\ell+H_\ell)$ contains $y_\ell$,
this implies \eqref{eq:fixed-parent-rank}.
Moreover, a positive-dimensional $D_\ell$ forces $H_\ell\neq0$.

By \cite[Proposition~6.1]{HabeggerPila2016}, the directions of the
geodesic closures of geodesic-optimal subvarieties of a fixed
$Y_\ell$ belong to a finite set. Take this set as $\mathcal H_\ell$.
The cited proposition applies to arbitrary subvarieties of complex
abelian varieties and imposes no condition on the translating points.
Thus the sets are independent of $y_\ell$, $Z_\ell$, $C_\ell$, and $g$.
\end{proof}

\subsection{Uniform positivity for the sequence}

We now return to sequence data over $\Qbar$. All abelian subgroup
directions produced by the preceding complex arguments are defined
over $\Qbar$. Indeed, an abelian subvariety of the complexification
of an abelian variety over $\Qbar$ is the image of a rational
endomorphism, by complete reducibility, and homomorphisms of abelian
varieties are unchanged by algebraically closed field extension.
Thus these directions are eligible for Lemma~\ref{lemma:projection}.

The finite sets in Lemma~\ref{lem:finite-directions} let us replace
directions depending on the marked points and subproducts by one
fixed direction along a subsequence. For a fixed product direction,
the tuple of local fiber dimensions takes only finitely many integer
values. Consequently
Lemma~\ref{lemma:projection} also excludes a fixed nonzero direction
whose sum of local fiber dimensions is at least its $\Phi$-rank for
infinitely many indices: pass to a subsequence on which that tuple is
constant. We use this equivalent form in the following proposition.

\begin{proposition}[Uniform tangent positivity]\label{prop:tangent}
Consider fixed sequence data as in Definition~\ref{def:sequence-data},
and let $V$ and $\Phi$ be as in Theorem~\ref{thm:sequence}.
Suppose that the conclusion of Lemma~\ref{lemma:projection} excludes
every fixed nonzero product $C=\prod_\ell C_\ell$ for which
\begin{equation}\label{eq:excluded-tangent-configuration}
 \sum_\ell\dim_{y_{i\ell}}
 \bigl(Y_\ell\cap(y_{i\ell}+C_\ell)\bigr)
 \geq\rank\bigl(\Phi|_{\Lie(C)}\bigr)
\end{equation}
holds for infinitely many $i$.
Put $r=\dim\cA+1$, and, for a linear map $g$ on $\Lie(\cA)$, define
\[
 \Delta g(v_1,\ldots,v_r)
   =\bigl(g(v_2)-g(v_1),\ldots,g(v_r)-g(v_{r-1})\bigr).
\]
Then the following assertions hold.
\begin{enumerate}
 \item There are a neighborhood $\mathcal U$ of $\Phi$ in
 $\End(\cA)_{\mathbb R}$ and an index $i_0$ such that
 $P_Z(g)>0$ for every $i\geq i_0$, every $g\in\mathcal U$, and every
 subproduct $Z=\prod_\ell Z_\ell$ with
 $y_{i\ell}\in Z_\ell\subseteq Y_\ell$.
 \item For every sufficiently large $i$ and every such subproduct $Z$,
 $P_{Z^r}(\Delta\Phi)>0$.
\end{enumerate}
\end{proposition}

\begin{proof}
Suppose first that the first assertion fails. Choosing successively
smaller neighborhoods and larger indices produces a subsequence,
relabeled by $i$, maps $g_i\longrightarrow\Phi$, and subproducts
$Z_i=\prod_\ell Z_{i\ell}$ through the marked points such that
$P_{Z_i}(g_i)=0$.
Lemma~\ref{lem:ax-product} gives product subgroups
$C_i=\prod_\ell C_{i\ell}$ satisfying
\[
 \sum_\ell\dim_{y_{i\ell}}
 \bigl(Z_{i\ell}\cap(y_{i\ell}+C_{i\ell})\bigr)
 \geq\rank\bigl(g_i|_{\Lie(C_i)}\bigr)+1.
\]
The sum on the left is positive. Apply
Lemma~\ref{lem:finite-directions} with $\varepsilon=1$ and pass to a
subsequence on which all the resulting directions $H_\ell$ are fixed.
Their product $H$ is nonzero. Passing again to a subsequence, we may
also fix the integers
\[
 a_\ell=\dim_{y_{i\ell}}
 \bigl(Y_\ell\cap(y_{i\ell}+H_\ell)\bigr).
\]
They are bounded between $0$ and $\dim Y_\ell$. We obtain
$\rank(g_i|_{\Lie(H)})\leq\sum_\ell a_\ell-1$.
Since a rank upper bound is closed under limits of linear maps,
\[
 \rank(\Phi|_{\Lie(H)})\leq\sum_\ell a_\ell-1.
\]
This contradicts the assumed exclusion of
\eqref{eq:excluded-tangent-configuration}.

For the second assertion, suppose that there are arbitrarily large
indices $i$ and subproducts $Z_i$ through the marked points with
$P_{Z_i^r}(\Delta\Phi)=0$. Apply Lemma~\ref{lem:ax-product} to the
product whose factors are indexed by both $\ell\in I$ and
$j\in\{1,\ldots,r\}$. For a fixed $i$, suppress $i$ from the notation,
write $C_j=\prod_\ell C_{j\ell}$ for the resulting subgroup in copy $j$,
and set
\[
 a_j=\sum_\ell\dim_{y_{i\ell}}
       \bigl(Z_{i\ell}\cap(y_{i\ell}+C_{j\ell})\bigr),
 \qquad E_j=\Phi(\Lie(C_j)).
\]
The consecutive-difference map on $\bigoplus_j E_j$ has kernel the
diagonal copy of $\bigcap_j E_j$. Its rank is consequently
$\sum_j\dim E_j-\dim\bigcap_j E_j$, and the lemma gives
\begin{equation}\label{eq:difference-direction-deficit}
 \sum_{j=1}^r a_j>
 \sum_{j=1}^r\dim E_j-\dim\bigcap_{j=1}^r E_j.
\end{equation}
There is an index $j$ such that $C_j\neq0$ and $a_j\geq\dim E_j$.
Indeed, if some $C_j=0$, then the intersection term vanishes and
\eqref{eq:difference-direction-deficit} forces a strict surplus in
another copy. If all $C_j$ are nonzero and every $a_j<\dim E_j$,
the integral deficits sum to at least
$r=\dim\cA+1$, whereas $\dim\bigcap_jE_j\leq\dim\cA$, again
contradicting \eqref{eq:difference-direction-deficit}.

Select one such copy for each $i$. If its $a_j$ is zero for any index,
then its nonzero product subgroup is contained in $\ker\Phi$.
That one fixed subgroup satisfies
\eqref{eq:excluded-tangent-configuration} for every index, since all
local fiber dimensions are nonnegative. This is impossible.
Thus the selected sum is positive. Apply
Lemma~\ref{lem:finite-directions} with $g=\Phi$ and $\varepsilon=0$.
It supplies a nonzero product direction in a fixed finite set and
satisfying the parent inequality. Passing to a fixed direction along
an infinite subsequence contradicts
\eqref{eq:excluded-tangent-configuration} once more.
\end{proof}

The first assertion supplies one neighborhood for all subproducts
through the tail of marked points. Section~\ref{sec:high} will choose
a rational box inside it and derive a degree-dependent numerical
lower bound for the tangent integrals. The second assertion gives
positivity at the limiting consecutive-difference map. In
Section~\ref{sec:arithmetic}, its numerical lower bound will follow
after the auxiliary products have fixed cycle classes. Neither
assertion here is a uniform positive numerical bound over all
subproducts.
